%----------------------------------------------------------------------------------------
% DE GONIOMETRISCHE VORM. Alle opgaven zijn eigen werk.
%
% Opbouw: Start, leerstof (theorie, werkwijze of uitleg, voorbeeld, korte oefening), en op een
% nieuwe pagina de Opdrachten (door elkaar, stijgende moeilijkheid), Koppel en Test jezelf.
% Labels van dit hoofdstuk beginnen met h3-.
%----------------------------------------------------------------------------------------

\mychapter{3}{De goniometrische vorm}

\lkr{Leerplan B+S'': LPD 41 (cartesische vorm omzetten naar goniometrische vorm en omgekeerd) en het eerste deel van LPD 42 (vermenigvuldigen en delen in goniometrische vorm), met de wenken over het verband met de somformules en de meetkundige betekenis. Machten en $n$-de machtswortels volgen in hoofdstuk 4. Het argument staat in radialen, met de hoofdwaarde in $]-\pi, \pi]$; wie liever in graden werkt, kan dat overal doen.}

In hoofdstuk 1 zag je dat vermenigvuldigen met $i$ een kwartslag is: het punt draait om de oorsprong, en zijn afstand tot de oorsprong blijft dezelfde. Dat is geen toeval voor $i$ alleen. Vermenigvuldigen met een complex getal is altijd een draaiing gevolgd door een uitrekking. Je ziet dat het best als je een complex getal niet beschrijft met zijn reëel en imaginair deel, maar met zijn afstand tot de oorsprong en een hoek. Dat is de goniometrische vorm.

\begin{oefeningen}[Start]
\oef{h3-start-1}

\oef{h3-verken}
\end{oefeningen}

\lkr{Laat oefening \ref{oef:h3-verken} maken vóór de uitleg. Leerlingen meten de hoeken in de figuur, in graden of in radialen. Het patroon (moduli vermenigvuldigen, hoeken optellen) is de kern van het hoofdstuk.}

%========================================================================================
\section{Modulus en argument}

Een punt in het vlak ligt vast door zijn twee coördinaten, maar ook door zijn afstand tot de oorsprong en de richting waarin het ligt. Voor een complex getal is die afstand de modulus. De richting geven we met een hoek, zoals op de goniometrische cirkel.

\subsection{Het argument}

\begin{theorie}
\begin{minipage}[t]{.57\linewidth}
\vspace{0pt}
Een \textbf{\textcolor{accent}{argument}} van een complex getal $z \neq 0$ is een hoek $\theta$ van de positieve reële as naar de vector $z$, tegen de wijzers van de klok in gemeten.

\vspace{1mm}
De \textbf{\textcolor{accent}{hoofdwaarde}} $\formule{\arg z}$ is het argument in $\formule{]-\pi, \pi]}$.
\end{minipage}\hfill
\begin{minipage}[t]{.38\linewidth}
\vspace{0pt}
\begin{tikzpicture}[x=.8cm,y=.8cm,font=\small]
  \draw[->,grijs] (-1,0) -- (4.2,0) node[right,text=black,font=\scriptsize] {Re};
  \draw[->,grijs] (0,-.8) -- (0,3.2) node[above,text=black,font=\scriptsize] {Im};
  \draw[vector] (0,0) -- node[above left=-1pt,text=black] {$r$} (3,2.4);
  \draw[accent,thick] (1,0) arc[start angle=0,end angle=38.66,radius=1];
  \node[text=black] at (19:1.35) {$\theta$};
  \cpunt{3}{2.4}{$z$}
\end{tikzpicture}
\end{minipage}
\end{theorie}

Een argument ligt niet helemaal vast: tel je er $2\pi$ bij, dan wijst de vector in dezelfde richting. Daarom spreken we een hoofdwaarde af. Getallen boven de reële as hebben een hoofdwaarde tussen 0 en $\pi$, getallen eronder een negatieve. Het getal 0 heeft geen argument, want de vector heeft geen richting.

\figuren{\begin{tikzpicture}[x=.55cm,y=.55cm,font=\scriptsize]
  \draw[->,grijs] (-3,0) -- (3.4,0) node[right,text=black] {Re};
  \draw[->,grijs] (0,-2.6) -- (0,2.9) node[above,text=black] {Im};
  \foreach \x/\y/\l/\p in {2.5/0/{0}/above,0/2.5/{\frac{\pi}{2}}/right,-2.5/0/{\pi}/above,0/-2.5/{-\frac{\pi}{2}}/right}{\cpunt[\p]{\x}{\y}{$\l$}}
  \node[below,align=center,text=black] at (0,-3) {argument van de getallen\\op de assen};
\end{tikzpicture}\qquad\qquad
\begin{tikzpicture}[x=.55cm,y=.55cm,font=\scriptsize]
  \fill[accent!12] (0,0) -- (3,0) arc[start angle=0,end angle=180,radius=3] -- cycle;
  \fill[fel!20] (0,0) -- (-3,0) arc[start angle=180,end angle=360,radius=3] -- cycle;
  \draw[->,grijs] (-3.2,0) -- (3.4,0) node[right,text=black] {Re};
  \draw[->,grijs] (0,-3.1) -- (0,3.3) node[above,text=black] {Im};
  \fill[white] (0,0) circle (2.5pt); \draw[grijs] (0,0) circle (2.5pt);
  \node[text=black,fill=accent!12,inner sep=1pt] at (0,1.6) {$0 < \arg z < \pi$};
  \node[text=black,fill=fel!20,inner sep=1pt] at (0,-1.6) {$-\pi < \arg z < 0$};
  \node[below,align=center,text=black] at (0,-3.2) {boven en onder de reële as\\(0 heeft geen argument)};
\end{tikzpicture}}

\voorbeeld{}{Geef de hoofdwaarde van het argument van $z = -2 + 2i$, en nog twee andere argumenten.\label{vb:h3-arg}}{%
\uitleg{Een schets toont waar het getal ligt. Ligt het op een as of op een bissectrice, dan lees je de hoek meteen af.}
\naast{\begin{minipage}{10cm}
\begin{stappenvrij}{3.6cm}{6.1cm}
\text{tweede kwadrant} & Reëel deel negatief, imaginair deel positief. \\[1ex]
\arg z = \pi - \tfrac{\pi}{4} = \tfrac{3\pi}{4} & Het punt ligt op de bissectrice: de hoek met de negatieve reële as is $\frac{\pi}{4}$. \\[1ex]
\tfrac{11\pi}{4} \ \text{en} \ {-\tfrac{5\pi}{4}} & Tel $2\pi$ bij $\frac{3\pi}{4}$, of trek $2\pi$ af: ook argumenten, maar geen hoofdwaarde.
\end{stappenvrij}
\end{minipage}}{\begin{tikzpicture}[x=.6cm,y=.6cm]
  \gauss{-3}{1}{-1}{3}
  \draw[vector] (0,0) -- (-2,2);
  \draw[accent,thick] (.7,0) arc[start angle=0,end angle=135,radius=.7];
  \cpunt[left]{-2}{2}{$z$}
\end{tikzpicture}}}

\begin{oefeningen}
\oef{h3-arg}
\end{oefeningen}

\subsection{De goniometrische vorm}

Ken je de modulus en een argument van een complex getal, dan ken je het getal: je weet hoe ver het ligt en in welke richting. Maar hoe vind je dan het reëel en het imaginair deel? Dat zegt de goniometrie.

\begin{theorie}
\begin{stelling}{Goniometrische vorm}
\begin{minipage}[t]{.6\linewidth}
\vspace{0pt}
Heeft $z = a + bi \neq 0$ modulus $r$ en argument $\theta$, dan is

\vspace{1mm}
\qquad $\formule{a = r\cos\theta}$ \ en \ $\formule{b = r\sin\theta}$, \ dus

\vspace{1mm}
\qquad $\formule{z = r(\cos\theta + i\sin\theta)}$
\end{minipage}\hfill
\begin{minipage}[t]{.36\linewidth}
\vspace{0pt}
\begin{tikzpicture}[x=.62cm,y=.62cm,font=\small]
  \draw[->,grijs] (-.5,0) -- (4,0) node[right,text=black,font=\scriptsize] {Re};
  \draw[->,grijs] (0,-.5) -- (0,3) node[above,text=black,font=\scriptsize] {Im};
  \draw[hulplijn] (1,0) arc[start angle=0,end angle=90,radius=1];
  \draw[hulplijn] (3,0) node[below,text=black,font=\scriptsize] {$a$} -- (3,2.4) -- (0,2.4) node[left,text=black,font=\scriptsize] {$b$};
  \draw[vector] (0,0) -- (3,2.4);
  \fill[accent] (.781,.625) circle (1.6pt);
  \node[left,font=\scriptsize] at (.75,.8) {$\frac{z}{r}$};
  \node[font=\scriptsize] at (19:1.45) {$\theta$};
  \cpunt{3}{2.4}{$z$}
\end{tikzpicture}
\end{minipage}
\end{stelling}

\begin{bewijs}
Deel het reëel en het imaginair deel door $r = \sqrt{a^2 + b^2}$. Het punt $\bigl(\frac{a}{r}, \frac{b}{r}\bigr)$ ligt op de goniometrische cirkel, want
\bewijsregel{\Bigl(\frac{a}{r}\Bigr)^2 + \Bigl(\frac{b}{r}\Bigr)^2 = \frac{a^2 + b^2}{r^2} = 1 ,}
en het ligt in dezelfde richting als $z$, dus bij de hoek $\theta$. Volgens de definitie van cosinus en sinus op de goniometrische cirkel is dan $\cos\theta = \frac{a}{r}$ en $\sin\theta = \frac{b}{r}$, of $a = r\cos\theta$ en $b = r\sin\theta$.
\end{bewijs}

De schrijfwijze $z = r(\cos\theta + i\sin\theta)$, met $r > 0$, is de \textbf{\textcolor{accent}{goniometrische vorm}} van $z$.
\end{theorie}

De vorm is pas goniometrisch als er precies $r(\cos\theta + i\sin\theta)$ staat, met $r > 0$ en twee keer dezelfde hoek. In $2(\cos\frac{\pi}{3} - i\sin\frac{\pi}{3})$ staat een minteken. Omdat $\cos(-\alpha) = \cos\alpha$ en $\sin(-\alpha) = -\sin\alpha$, is dat $2(\cos(-\frac{\pi}{3}) + i\sin(-\frac{\pi}{3}))$, met argument $-\frac{\pi}{3}$. Ook een negatieve factor voor de haakjes kun je wegwerken, met de formules $\cos(\alpha + \pi) = -\cos\alpha$ en $\sin(\alpha + \pi) = -\sin\alpha$.

\voorbeeld{}{Schrijf $z = -2\bigl(\cos\frac{\pi}{3} + i\sin\frac{\pi}{3}\bigr)$ in goniometrische vorm, met de hoofdwaarde van het argument.\label{vb:h3-geengonio}}{%
\uitleg{De factor $-2$ is negatief, dus dit is geen goniometrische vorm. Het minteken brengen we in de haakjes, en we zoeken een hoek waarbij cosinus en sinus van teken veranderen: een halve toer verder.}
\begin{stappenbreed}
z = 2\Bigl(-\cos\frac{\pi}{3} - i\sin\frac{\pi}{3}\Bigr) & Het minteken naar binnen. \\[3ex]
= 2\Bigl(\cos\frac{4\pi}{3} + i\sin\frac{4\pi}{3}\Bigr) & $\cos(\frac{\pi}{3} + \pi) = -\cos\frac{\pi}{3}$, en zo ook voor de sinus. \\[3ex]
= 2\Bigl(\cos\Bigl(-\frac{2\pi}{3}\Bigr) + i\sin\Bigl(-\frac{2\pi}{3}\Bigr)\Bigr) & De hoofdwaarde: $\frac{4\pi}{3} - 2\pi = -\frac{2\pi}{3}$.
\end{stappenbreed}
Controle in cartesische vorm: $-2\bigl(\frac12 + \frac{\sqrt3}{2}i\bigr) = -1 - i\sqrt3$, en $2\bigl(-\frac12 - \frac{\sqrt3}{2}i\bigr) = -1 - i\sqrt3$.}

\begin{oefeningen}
\oef{h3-gonio}
\end{oefeningen}

\subsection{Gelijke getallen in goniometrische vorm}

In het vorige voorbeeld gaf de hoek $\frac{4\pi}{3}$ hetzelfde getal als $-\frac{2\pi}{3}$. Wanneer stellen twee goniometrische vormen hetzelfde getal voor? Je kunt dat zien zonder naar de cartesische vorm te gaan.

\begin{theorie}
\begin{stelling}{Gelijke getallen in goniometrische vorm}
Voor $r_1, r_2 > 0$ is \quad $\formule{r_1(\cos\theta_1 + i\sin\theta_1) = r_2(\cos\theta_2 + i\sin\theta_2)}$

\vspace{1mm}
als en slechts als \quad $\formule{r_1 = r_2}$ \ en \ $\formule{\theta_1 = \theta_2 + k \cdot 2\pi}$ \ voor een geheel getal $k$.
\end{stelling}

\begin{bewijs}
Zijn de moduli gelijk en verschillen de hoeken $k \cdot 2\pi$, dan zijn de cosinussen en de sinussen gelijk, want cosinus en sinus hebben periode $2\pi$. Stel omgekeerd dat de getallen gelijk zijn. Dan hebben ze dezelfde modulus, en
\bewijsregel{|r(\cos\theta + i\sin\theta)| = \sqrt{r^2\cos^2\theta + r^2\sin^2\theta} = \sqrt{r^2} = r ,}
dus $r_1 = r_2$. Deel door die modulus: $\cos\theta_1 + i\sin\theta_1 = \cos\theta_2 + i\sin\theta_2$. Gelijke complexe getallen hebben gelijke reële en imaginaire delen, dus $\cos\theta_1 = \cos\theta_2$ en $\sin\theta_1 = \sin\theta_2$. De hoeken $\theta_1$ en $\theta_2$ horen dus bij hetzelfde punt van de goniometrische cirkel, en verschillen een veelvoud van $2\pi$.
\end{bewijs}
\end{theorie}

Zo breng je elk argument terug naar de hoofdwaarde: tel er een veelvoud van $2\pi$ bij of trek het ervan af, tot de hoek in $]-\pi, \pi]$ ligt. In hoofdstuk 4 heb je deze stelling nodig om machtswortels te zoeken.

\voorbeeld{}{Schrijf $z = 3\bigl(\cos\frac{11\pi}{4} + i\sin\frac{11\pi}{4}\bigr)$ met de hoofdwaarde van het argument.\label{vb:h3-gelijk}}{%
\uitleg{De hoek $\frac{11\pi}{4}$ is groter dan $\pi$. We trekken er $2\pi = \frac{8\pi}{4}$ van af, tot hij in $]-\pi, \pi]$ ligt. De modulus blijft dezelfde.}
\begin{stappen}
\frac{11\pi}{4} - 2\pi = \frac{11\pi - 8\pi}{4} = \frac{3\pi}{4} & Een volle toer minder. \\[3ex]
z = 3\Bigl(\cos\frac{3\pi}{4} + i\sin\frac{3\pi}{4}\Bigr) & $\frac{3\pi}{4}$ ligt in $]-\pi, \pi]$.
\end{stappen}}

\begin{oefeningen}
\oef{h3-hoofdwaarde}
\end{oefeningen}

%========================================================================================
\section{Omzetten}

\subsection{Van cartesisch naar goniometrisch}

De modulus ken je al. Voor het argument gebruik je dat $\cos\theta = \frac{a}{r}$ en $\sin\theta = \frac{b}{r}$. Bij een bijzondere hoek herken je die waarden op de goniometrische cirkel.

\voorbeeld{}{Schrijf $z = -1 + i\sqrt3$ in goniometrische vorm.\label{vb:h3-omzet}}{%
\uitleg{We berekenen de modulus, en zoeken dan een hoek met de juiste cosinus en sinus. Een schets toont in welk kwadrant $z$ ligt.}
\naast{\begin{minipage}{10.4cm}
\begin{stappenvrij}{4.6cm}{5.5cm}
r = \sqrt{1 + 3} = 2 & De modulus. \\[1ex]
\cos\theta = -\tfrac12, \quad \sin\theta = \tfrac{\sqrt3}{2} & Deel $a$ en $b$ door $r$. \\[1ex]
\theta = \tfrac{2\pi}{3} & Het tweede kwadrant: $\pi - \tfrac{\pi}{3}$. \\[1ex]
z = 2\bigl(\cos\tfrac{2\pi}{3} + i\sin\tfrac{2\pi}{3}\bigr) & De goniometrische vorm.
\end{stappenvrij}
\end{minipage}}{\begin{tikzpicture}[x=.6cm,y=.6cm,font=\small]
  \gauss{-2}{2}{-1}{2}
  \draw[vector] (0,0) -- (-1,1.732);
  \draw[accent,thick] (.6,0) arc[start angle=0,end angle=120,radius=.6];
  \cpunt[above left]{-1}{1.732}{$z$}
\end{tikzpicture}}}

\voorbeeld{}{Schrijf $z = -3 - 3i$ in goniometrische vorm.\label{vb:h3-omzet2}}{%
\uitleg{Het getal ligt in het derde kwadrant. We berekenen de modulus en zoeken een hoek in $]-\pi, -\frac{\pi}{2}[$.}
\naast{\begin{minipage}{10.9cm}
\begin{stappenvrij}{6.1cm}{4.5cm}
r = \sqrt{9 + 9} = 3\sqrt2 & De modulus. \\[1ex]
\cos\theta = \sin\theta = \frac{-3}{3\sqrt2} = -\frac{\sqrt2}{2} & Een bijzondere hoek. \\[3ex]
\theta = -\tfrac{3\pi}{4} & In het derde kwadrant, met een negatieve hoofdwaarde: $-\pi + \frac{\pi}{4}$. \\[1ex]
z = 3\sqrt2\bigl(\cos(-\tfrac{3\pi}{4}) + i\sin(-\tfrac{3\pi}{4})\bigr) &
\end{stappenvrij}
\end{minipage}}{\begin{tikzpicture}[x=.55cm,y=.55cm,font=\small]
  \gauss{-3}{1}{-3}{1}
  \draw[vector] (0,0) -- (-3,-3);
  \draw[accent,thick] (.6,0) arc[start angle=0,end angle=-135,radius=.6];
  \cpunt[left]{-3}{-3}{$z$}
\end{tikzpicture}}}

Ligt de hoek niet op de goniometrische cirkel, dan gebruik je je rekentoestel. In het eerste en het vierde kwadrant geeft de tangens meteen de juiste hoek.

\voorbeeld{}{Schrijf $z = 4 - 3i$ in goniometrische vorm, met het argument op drie decimalen.\label{vb:h3-omzet3}}{%
\uitleg{Dit is geen bijzondere hoek. Het getal ligt in het vierde kwadrant, en daar geeft de tangens meteen de juiste hoek.}
\begin{stappen}
r = \sqrt{16 + 9} = 5 & De modulus. \\[1ex]
\theta = \bgtan\bigl(-\tfrac34\bigr) \approx -0{,}644 & Vierde kwadrant: de rekentoestel geeft de juiste hoek, in radialen. \\[1ex]
z \approx 5\bigl(\cos(-0{,}644) + i\sin(-0{,}644)\bigr) &
\end{stappen}}

In het tweede en het derde kwadrant werkt de tangens niet zomaar. De toets $\tan^{-1}$ geeft altijd een hoek tussen $-\frac{\pi}{2}$ en $\frac{\pi}{2}$, dus in het eerste of het vierde kwadrant. En $\tan\theta = \frac{b}{a}$ maakt geen onderscheid tussen tegenovergestelde punten: voor $-1 - i$ en voor $1 + i$ is $\frac{b}{a}$ allebei 1. Wie bij $-1 - i$ zomaar $\bgtan 1 = \frac{\pi}{4}$ neemt, wijst de verkeerde kant op.

\voorbeeld{}{Schrijf $z = -5 + 2i$ in goniometrische vorm, met de modulus en het argument op drie decimalen.\label{vb:h3-omzet4}}{%
\uitleg{Het getal ligt in het tweede kwadrant. De rekentoestel geeft een hoek in het vierde kwadrant, de richting van het tegengestelde punt. We draaien die nog een halve toer.}
\naast{\begin{minipage}{9.8cm}
\begin{stappenvrij}{5.3cm}{4.2cm}
r = \sqrt{25 + 4} = \sqrt{29} \approx 5{,}385 & De modulus. \\[1ex]
\bgtan\bigl(-\tfrac25\bigr) \approx -0{,}381 & Een hoek in het vierde kwadrant: de richting van $5 - 2i$. \\[1ex]
\theta \approx -0{,}381 + \pi \approx 2{,}761 & Een halve toer verder: het tweede kwadrant. \\[1ex]
z \approx 5{,}385(\cos 2{,}761 + i\sin 2{,}761) &
\end{stappenvrij}
\end{minipage}}{\begin{tikzpicture}[x=.4cm,y=.4cm,font=\small]
  \gaussstap{2}{-5}{5}{-2}{2}
  \draw[hulplijn] (-5,2) -- (5,-2);
  \draw[vector] (0,0) -- (-5,2);
  \draw[accent,thick] (.8,0) arc[start angle=0,end angle=158.2,radius=.8];
  \draw[grijs,thick] (1.6,0) arc[start angle=0,end angle=-21.8,radius=1.6];
  \cpunt[above]{-5}{2}{$z$}
  \fill[grijs] (5,-2) circle (2pt) node[below,text=black,font=\scriptsize] {$5 - 2i$};
\end{tikzpicture}}}

\begin{werkwijze}
Schets het punt, zodat je het kwadrant kent. Bereken $r = \sqrt{a^2 + b^2}$. Zoek $\theta$ met $\cos\theta = \frac{a}{r}$ en $\sin\theta = \frac{b}{r}$, of met $\bgtan\frac{b}{a}$: in het tweede kwadrant tel je er dan $\pi$ bij, in het derde trek je er $\pi$ van af.
\end{werkwijze}

\nujij{Schrijf $1 - i$ in goniometrische vorm.}{$r = \sqrt2$ en $\theta = -\frac{\pi}{4}$: $\sqrt2\bigl(\cos(-\frac{\pi}{4}) + i\sin(-\frac{\pi}{4})\bigr)$.}

\begin{oefeningen}
\oef{h3-omzet}

\oef{h3-rekentoestel}
\end{oefeningen}

\subsection{Van goniometrisch naar cartesisch}

De omgekeerde weg is korter: je rekent de cosinus en de sinus uit, en werkt de haakjes uit.

\voorbeeld{}{Schrijf $z = 4\bigl(\cos\frac{5\pi}{6} + i\sin\frac{5\pi}{6}\bigr)$ in cartesische vorm.\label{vb:h3-terug}}{%
\uitleg{De hoek $\frac{5\pi}{6}$ ligt in het tweede kwadrant: de cosinus is negatief en de sinus positief.}
\begin{stappen}
z = 4\Bigl(-\frac{\sqrt3}{2} + \frac12 i\Bigr) & $\cos\frac{5\pi}{6} = -\frac{\sqrt3}{2}$ en $\sin\frac{5\pi}{6} = \frac12$. \\[3ex]
z = -2\sqrt3 + 2i & Haakjes uitwerken.
\end{stappen}}

\begin{werkwijze}
Bereken $\cos\theta$ en $\sin\theta$, exact bij een bijzondere hoek en anders met het rekentoestel in radialen. Vermenigvuldig ze met $r$: dat zijn het reëel en het imaginair deel.
\end{werkwijze}

\begin{oefeningen}
\oef{h3-terug}
\end{oefeningen}

%========================================================================================
\section{Vermenigvuldigen en delen}

\subsection{Het product}

In de Verken werd de modulus bij elke vermenigvuldiging met $1 + i$ met $\sqrt2$ vermenigvuldigd, en kwam er bij de hoek telkens $\frac{\pi}{4}$ bij. Dat geldt voor elk product, en het bewijs gebruikt de somformules uit de goniometrie.

\begin{theorie}
\begin{stelling}{Product in goniometrische vorm}
Voor $z_1 = r_1(\cos\theta_1 + i\sin\theta_1)$ en $z_2 = r_2(\cos\theta_2 + i\sin\theta_2)$ is

\vspace{1mm}
\qquad $\formule{z_1 z_2 = r_1 r_2 \bigl(\cos(\theta_1 + \theta_2) + i\sin(\theta_1 + \theta_2)\bigr)}$
\end{stelling}

\begin{bewijs}
Werk het product uit, en vervang $i^2$ door $-1$:
\bewijsregel{z_1 z_2 = r_1 r_2 \bigl((\cos\theta_1\cos\theta_2 - \sin\theta_1\sin\theta_2) + i(\sin\theta_1\cos\theta_2 + \cos\theta_1\sin\theta_2)\bigr) .}
Volgens de somformules is het eerste haakje $\cos(\theta_1 + \theta_2)$ en het tweede $\sin(\theta_1 + \theta_2)$.
\end{bewijs}
\end{theorie}

\naast{\begin{minipage}{9.6cm}
Bij een product vermenigvuldig je de moduli en tel je de argumenten op. Meetkundig: vermenigvuldigen met $z_2$ draait $z_1$ over de hoek $\theta_2$ om de oorsprong, en rekt de afstand tot de oorsprong uit met de factor $r_2$. Zo is vermenigvuldigen met $i = \cos\frac{\pi}{2} + i\sin\frac{\pi}{2}$ een kwartslag zonder uitrekking: dat zag je al in hoofdstuk 1.

Wat in cartesische vorm een hele uitwerking was, is hier een vermenigvuldiging en een optelling.
\end{minipage}}{\begin{tikzpicture}[x=.62cm,y=.62cm,font=\small]
  \draw[->,grijs] (-3.3,0) -- (3.4,0) node[right,text=black,font=\scriptsize] {Re};
  \draw[->,grijs] (0,-.4) -- (0,4.4) node[above,text=black,font=\scriptsize] {Im};
  \draw[vector] (0,0) -- (2.6,1.5);
  \draw[vector,grijs] (0,0) -- (.345,1.288);
  \draw[vector,fel!80!black] (0,0) -- (-1.04,3.86);
  \draw[accent,thick] (1,0) arc[start angle=0,end angle=30,radius=1];
  \draw[fel!80!black,thick] (30:2.3) arc[start angle=30,end angle=105,radius=2.3];
  \node at (15:1.35) {$\theta_1$};
  \node at (66:2.7) {$\theta_2$};
  \cpunt[right]{2.6}{1.5}{$z_1$}
  \cpunt[right]{.345}{1.288}{$z_2$}
  \cpunt[left]{-1.04}{3.86}{$z_1 z_2$}
\end{tikzpicture}}

\voorbeeld{}{Bereken $z_1 z_2$ voor $z_1 = 2\bigl(\cos\frac{\pi}{3} + i\sin\frac{\pi}{3}\bigr)$ en $z_2 = 3\bigl(\cos\frac{\pi}{4} + i\sin\frac{\pi}{4}\bigr)$.\label{vb:h3-product}}{%
\uitleg{Beide getallen staan in goniometrische vorm: we vermenigvuldigen de moduli en tellen de argumenten op.}
\begin{stappen}
r_1 r_2 = 2 \cdot 3 = 6 & De moduli vermenigvuldigen. \\[1ex]
\theta_1 + \theta_2 = \frac{\pi}{3} + \frac{\pi}{4} = \frac{7\pi}{12} & De argumenten optellen. \\[3ex]
z_1 z_2 = 6\Bigl(\cos\frac{7\pi}{12} + i\sin\frac{7\pi}{12}\Bigr) & In de cartesische vorm zou je $\cos\frac{7\pi}{12}$ nodig hebben, en dat is geen bijzondere waarde.
\end{stappen}}

De som van twee hoofdwaarden ligt niet altijd in $]-\pi, \pi]$. Valt ze erbuiten, tel er dan $2\pi$ bij of trek er $2\pi$ van af.

\voorbeeld{}{Bepaal de modulus en de hoofdwaarde van het argument van $z_1 z_2$, voor $z_1 = 2\bigl(\cos\frac{3\pi}{4} + i\sin\frac{3\pi}{4}\bigr)$ en $z_2 = 3\bigl(\cos\frac{2\pi}{3} + i\sin\frac{2\pi}{3}\bigr)$.\label{vb:h3-product-terug}}{%
\uitleg{We vermenigvuldigen de moduli en tellen de argumenten op. Twee hoeken in het tweede kwadrant geven samen meer dan $\pi$: we brengen de som terug.}
\begin{stappen}
|z_1 z_2| = 2 \cdot 3 = 6 & De moduli vermenigvuldigen. \\[1ex]
\frac{3\pi}{4} + \frac{2\pi}{3} = \frac{9\pi + 8\pi}{12} = \frac{17\pi}{12} & Groter dan $\pi = \frac{12\pi}{12}$. \\[3ex]
\arg(z_1 z_2) = \frac{17\pi}{12} - 2\pi = -\frac{7\pi}{12} & Een volle toer minder. Het product ligt onder de reële as.
\end{stappen}}

Voor een som bestaat zo'n regel niet. Tel je $\cos 0 + i\sin 0 = 1$ en $\cos\frac{\pi}{2} + i\sin\frac{\pi}{2} = i$ op, dan vind je $1 + i$, met modulus $\sqrt2$ en niet $1 + 1 = 2$. Optellen en aftrekken doe je dus in cartesische vorm, vermenigvuldigen en delen vaak in goniometrische vorm.

\subsection{Het quotiënt}

Delen is het omgekeerde van vermenigvuldigen. Dus deel je de moduli, en trek je de argumenten af.

\begin{theorie}
\begin{stelling}{Quotiënt in goniometrische vorm}
Voor $z_1 = r_1(\cos\theta_1 + i\sin\theta_1)$ en $z_2 = r_2(\cos\theta_2 + i\sin\theta_2)$, met $z_2 \neq 0$, is

\vspace{1mm}
\qquad $\formule{\dfrac{z_1}{z_2} = \dfrac{r_1}{r_2} \bigl(\cos(\theta_1 - \theta_2) + i\sin(\theta_1 - \theta_2)\bigr)}$
\end{stelling}

\begin{bewijs}
Noem $w = \frac{r_1}{r_2}\bigl(\cos(\theta_1 - \theta_2) + i\sin(\theta_1 - \theta_2)\bigr)$. Met de regel voor het product is
\bewijsregel{w \cdot z_2 = \frac{r_1}{r_2} \cdot r_2 \bigl(\cos(\theta_1 - \theta_2 + \theta_2) + i\sin(\theta_1 - \theta_2 + \theta_2)\bigr) = z_1 .}
Het quotiënt $\frac{z_1}{z_2}$ is het enige getal dat met $z_2$ vermenigvuldigd $z_1$ geeft, dus $w = \frac{z_1}{z_2}$.
\end{bewijs}
\end{theorie}

Neem je $z_1 = 1 = \cos 0 + i\sin 0$, dan vind je het omgekeerde: $\frac{1}{z_2} = \frac{1}{r_2}\bigl(\cos(-\theta_2) + i\sin(-\theta_2)\bigr)$. Het omgekeerde van een getal ligt dus in de richting van zijn toegevoegde, op afstand $\frac{1}{r_2}$.

\voorbeeld{}{Bereken $\dfrac{z_1}{z_2}$ voor dezelfde $z_1 = 2\bigl(\cos\frac{\pi}{3} + i\sin\frac{\pi}{3}\bigr)$ en $z_2 = 3\bigl(\cos\frac{\pi}{4} + i\sin\frac{\pi}{4}\bigr)$.\label{vb:h3-quotient}}{%
\uitleg{We delen de moduli en trekken de argumenten af, in de volgorde van het quotiënt.}
\begin{stappen}
\frac{z_1}{z_2} = \frac23\Bigl(\cos\Bigl(\frac{\pi}{3} - \frac{\pi}{4}\Bigr) + i\sin\Bigl(\frac{\pi}{3} - \frac{\pi}{4}\Bigr)\Bigr) & Moduli delen, argumenten aftrekken. \\[3ex]
= \frac23\Bigl(\cos\frac{\pi}{12} + i\sin\frac{\pi}{12}\Bigr) & $\frac{4\pi}{12} - \frac{3\pi}{12} = \frac{\pi}{12}$.
\end{stappen}}

\begin{werkwijze}
Zet de getallen in goniometrische vorm, als ze dat nog niet zijn. Bij een product vermenigvuldig je de moduli en tel je de argumenten op, bij een quotiënt deel je de moduli en trek je de argumenten af. Breng het argument terug naar $]-\pi, \pi]$, en zet het resultaat in cartesische vorm als de hoek bijzonder is.
\end{werkwijze}

\nujij{Bereken $zw$ en $\frac{z}{w}$ voor $z = 4\bigl(\cos\frac{3\pi}{4} + i\sin\frac{3\pi}{4}\bigr)$ en $w = 2\bigl(\cos\frac{\pi}{2} + i\sin\frac{\pi}{2}\bigr)$.}{$zw = 8\bigl(\cos\frac{5\pi}{4} + i\sin\frac{5\pi}{4}\bigr) = 8\bigl(\cos(-\frac{3\pi}{4}) + i\sin(-\frac{3\pi}{4})\bigr)$ en $\frac{z}{w} = 2\bigl(\cos\frac{\pi}{4} + i\sin\frac{\pi}{4}\bigr) = \sqrt2 + i\sqrt2$.}

\begin{oefeningen}
\oef{h3-product}

\oef{h3-product2}
\end{oefeningen}

\subsection{Draaien om de oorsprong}

Een getal met modulus 1 rekt niets uit. Vermenigvuldigen ermee is dus een zuivere draaiing. Zo kun je een punt draaien zonder goniometrie op de figuur zelf.

\begin{theorie}
\begin{stelling}{Draaiing om de oorsprong}
Draai je het punt $z$ om de oorsprong over de hoek $\alpha$, dan kom je in het punt

\vspace{1mm}
\qquad $\formule{z \cdot (\cos\alpha + i\sin\alpha)}$
\end{stelling}

\begin{bewijs}
Is $z = 0$, dan blijft het punt op zijn plaats, en is ook het product 0. Is $z = r(\cos\theta + i\sin\theta) \neq 0$, dan is volgens de regel voor het product
\bewijsregel{z \cdot (\cos\alpha + i\sin\alpha) = r\bigl(\cos(\theta + \alpha) + i\sin(\theta + \alpha)\bigr) .}
Dat punt ligt even ver van de oorsprong als $z$, en zijn argument is $\alpha$ groter: het is het beeld van $z$ bij de draaiing.
\end{bewijs}
\end{theorie}

Met $\alpha = \frac{\pi}{2}$ vind je de kwartslag van hoofdstuk 1: $\cos\frac{\pi}{2} + i\sin\frac{\pi}{2} = i$. Een negatieve hoek draait met de wijzers van de klok mee. En vermenigvuldigen met een willekeurig getal $\rho(\cos\alpha + i\sin\alpha)$ is een draaiing over $\alpha$, gevolgd door een uitrekking met factor $\rho$ vanuit de oorsprong.

\voorbeeld{}{Draai het punt $z = 2\sqrt3 + 2i$ over $\frac{\pi}{3}$ om de oorsprong.\label{vb:h3-draai}}{%
\uitleg{We vermenigvuldigen met $\cos\frac{\pi}{3} + i\sin\frac{\pi}{3}$. Dat kan in goniometrische vorm of in cartesische vorm; we doen het op beide manieren.}
\naast{\begin{minipage}{10.6cm}
\begin{stappenvrij}{5.4cm}{4.8cm}
z = 4\bigl(\cos\tfrac{\pi}{6} + i\sin\tfrac{\pi}{6}\bigr) & $r = \sqrt{12 + 4} = 4$ en $\theta = \frac{\pi}{6}$. \\[1ex]
4\bigl(\cos\tfrac{\pi}{2} + i\sin\tfrac{\pi}{2}\bigr) = 4i & Argument $\frac{\pi}{6} + \frac{\pi}{3} = \frac{\pi}{2}$. \\[1ex]
(2\sqrt3 + 2i)\Bigl(\frac12 + \frac{\sqrt3}{2} i\Bigr) & Cartesisch: $\cos\frac{\pi}{3} + i\sin\frac{\pi}{3} = \frac12 + \frac{\sqrt3}{2}i$. \\[3ex]
= \sqrt3 + 3i + i - \sqrt3 = 4i & Dezelfde uitkomst.
\end{stappenvrij}
\end{minipage}}{\begin{tikzpicture}[x=.55cm,y=.55cm,font=\small]
  \gaussstap{2}{-1}{4}{-1}{4}
  \draw[vector] (0,0) -- (3.464,2);
  \draw[vector,fel!80!black] (0,0) -- (0,4);
  \draw[hulplijn] (30:4) arc[start angle=30,end angle=90,radius=4];
  \cpunt[right]{3.464}{2}{$z$}
  \cpunt[right]{0}{4}{$4i$}
\end{tikzpicture}}}

\begin{werkwijze}
Vermenigvuldig het punt met $\cos\alpha + i\sin\alpha$, met $\alpha > 0$ voor een draaiing tegen de wijzers van de klok in en $\alpha < 0$ voor een draaiing met de wijzers mee. Reken in cartesische vorm als je het resultaat in cartesische vorm nodig hebt.
\end{werkwijze}

\begin{oefeningen}
\oef{h3-draai}
\end{oefeningen}

%========================================================================================
\opdrachtenkop

\begin{oefeningen}
\oef{h3-o4}

\oef{h3-o1}

\oef{h3-o16}

\oef{h3-o2}

\oef{h3-o15}

\oef{h3-o5}

\oef{h3-o17}

\oef{h3-o19}

\oef{h3-o3}

\oef{h3-o9}

\oef{h3-o18}

\oef{h3-o7}

\oef{h3-o22}

\oef{h3-o11}

\oef{h3-o6}

\oef{h3-o23}

\oef{h3-o8}

\oef{h3-o20}

\oef{h3-o21}

\oef{h3-o10}

\oef{h3-o12}

\oef{h3-o24}

\oef{h3-o14}

\oef{h3-o25}

\oef{h3-o13}

\oef{h3-o26}
\end{oefeningen}

%----------------------------------------------------------------------------------------
% KOPPEL: verbanden met andere leerstof
%----------------------------------------------------------------------------------------
\needspace{20\baselineskip}
\opdrachtsectie{Koppel}

\lkr{De somformules komen uit de cursus Goniometrische functies (5de jaar). Kijk na of ze al gezien zijn: het product in goniometrische vorm steunt erop. De tweede oefening bouwt voort op de eerste Koppel-oefening van hoofdstuk 1.}

\begin{oefeningen}
\oef{h3-koppel-1}

\oef{h3-koppel-2}
\end{oefeningen}

%----------------------------------------------------------------------------------------
\opdrachtsectie{Test jezelf}

\begin{oefeningen}
\oef{h3-t1}

\oef{h3-t2}

\oef{h3-t3}

\oef{h3-t4}

\oef{h3-t5}

\oef{h3-t6}

\oef{h3-t7}
\end{oefeningen}

\herhaalschema{%
  \fv{oef:h3-t1}{vb:h3-omzet}{oef:h3-omzet}%
  \fv{oef:h3-t2}{vb:h3-omzet2}{oef:h3-omzet}%
  \fv{oef:h3-t3}{vb:h3-terug}{oef:h3-terug}%
  \fv{oef:h3-t4}{vb:h3-product}{oef:h3-product}%
  \fv{oef:h3-t5}{vb:h3-draai}{oef:h3-draai}%
  \fv{oef:h3-t6}{vb:h3-omzet4}{oef:h3-rekentoestel}%
  \fv{oef:h3-t7}{vb:h3-product-terug}{oef:h3-product2}}

\ifzelfstudie
\section*{Maak je eigen samenvatting}
Sluit dit boek en schrijf in je schrift: wat het argument en de hoofdwaarde zijn, de goniometrische vorm, hoe je omzet van cartesisch naar goniometrisch en terug (en waar je moet opletten met de tangens), hoe je vermenigvuldigt en deelt in goniometrische vorm, en wat vermenigvuldigen meetkundig betekent. Vergelijk daarna met de kaders.
\fi
